\chapter{Machine Learning Background}

\section{Supervised Learning for Regression}

The pricing task is formulated as a supervised regression problem.
Each observation consists of option and market parameters, and the target is the corresponding Black-Scholes call price.
Given a dataset $\{(x_i, y_i)\}_{i=1}^{n}$, the goal is to learn a function $f_\theta$ such that:

\begin{equation}
    f_\theta(x_i) \approx y_i.
\end{equation}

In this project, the target $y_i$ is not observed from markets.
It is generated by the analytical Black-Scholes formula, which allows direct measurement of approximation error.

\section{Neural Networks as Function Approximators}

Feed-forward neural networks are flexible nonlinear function approximators.
Under general assumptions, sufficiently large feed-forward networks can approximate continuous functions on compact domains \cite{cybenko1989approximation,hornik1989multilayer}.
This makes them natural candidates for approximating a deterministic pricing function such as Black-Scholes.

The purpose of the neural network is not to replace the analytical formula when that formula is available.
Rather, the Black-Scholes setting provides a clean benchmark for studying the surrogate-model idea before moving to models where closed-form prices may not exist.

\section{Feed-Forward Neural Network Architecture}

The model used in the project is a Multi-Layer Perceptron.
It receives numerical input features and returns one scalar output, the predicted option price.

The baseline architecture is:

\begin{equation}
    5 \rightarrow 64 \rightarrow 64 \rightarrow 32 \rightarrow 1,
\end{equation}

where the five inputs are $(S_0, K, T, r, \sigma)$.
After intermediate experiments, the selected final configuration uses six input features, including moneyness, and SiLU hidden activations:

\begin{equation}
    6 \rightarrow 64 \rightarrow 64 \rightarrow 32 \rightarrow 1.
\end{equation}

The final neural network approximates:

\begin{equation}
    f_\theta(S_0, K, T, r, \sigma, S_0/K) \approx C_{BS}.
\end{equation}

\section{Support Vector Regression}

Support Vector Regression is included as an additional classical machine learning baseline.
The project uses SVR with an RBF kernel, which is able to model nonlinear relationships in tabular data.

SVR is evaluated on reduced datasets because kernel methods generally scale less favorably than neural networks as the number of training observations increases.
For this reason, the SVR experiments should be interpreted as a classical reference point rather than as a full-scale replacement for the neural network pipeline.

\section{Pricing Surrogate Models}

A pricing surrogate is a model trained to approximate a pricing function.
Once trained, the surrogate can produce prices quickly for new inputs.
This is useful when the original pricing method is computationally expensive, for example when repeated Monte Carlo simulations are required.

In the present Black-Scholes setting, the analytical formula remains the fastest and most exact method.
The relevant surrogate comparison is therefore mainly between neural network inference and Monte Carlo simulation.
